\subsection{Variational energies and feasible spaces}
\label{sec:models}

The three \(p\)-energies are defined in \cref{eq:pure-energy,eq:reg-energy,eq:reaction-energy}. The
remaining models have the following energies, written for \(d\leq2\):
\begin{align}
 \E_{\rm lin}(u)&=\int_\Omega\left(\frac12\abs{\nabla u}^2+
 \frac12u^2\right)\dd x-\ip fu,
 \label{eq:linear-energy}\\
 \E_{\rm cub}(u)&=\int_\Omega\left(\frac12\abs{\nabla u}^2+
 \frac14u^4\right)\dd x-\ip fu,
 \label{eq:cubic-energy}\\
 \E_{\sinh}(u)&=\int_\Omega\left(\frac12\abs{\nabla u}^2+
 \cosh u-1\right)\dd x-\ip fu.
 \label{eq:sinh-energy}
\end{align}
The pairing is the energy-space duality; for a regular load it is the corresponding integral.
The additive \(-1\) in \cref{eq:sinh-energy} normalizes the
potential and does not change the Euler--Lagrange equation.

\begin{assumption}[Data and feasible spaces for the six models]
\label[assumption]{ass:model-data}
The following model-specific conditions are imposed.
\begin{enumerate}[label=\textup{(\roman*)}]
 \item For the linear and cubic models, \(f\in (H^1(\Omega))^*\). For the sinh model, the convex
 integral functional has a load \(f\in(H^1(\Omega))^*\), and the
 integral functional is allowed to take the value \(+\infty\), its effective domain is nonempty, and
 local derivative/Hessian statements are restricted to an \(L^\infty\)-bounded neighborhood.
 \item For the pure and regularized \(p\)-energies, \(f\in X_{p,\circ}^*\), or equivalently the Neumann
 load annihilates constants before restriction to \(X_{p,\circ}\).
 \item For the reaction \(p\)-energy, \(f\in(W^{1,p})^*\). In the regularized model, the parameter
 \(\eps\) in \cref{eq:reg-energy} is fixed and positive.
\end{enumerate}
\end{assumption}

\begin{proposition}[Well-posedness]
\label{prop:wellposedness}
Under \Cref{ass:model-data}, \cref{eq:linear-energy,eq:cubic-energy,eq:sinh-energy} have unique minimizers
in \(H^1(\Omega)\); \cref{eq:pure-energy,eq:reg-energy} have unique minimizers in \(X_{p,\circ}\); and
\cref{eq:reaction-energy} has a unique minimizer in \(W^{1,p}(\Omega)\).
\end{proposition}

\begin{proof}[Proof sketch]
The quadratic, strictly convex, and $p$-growth terms give coercivity on the stated feasible spaces; strict convexity gives uniqueness.  The pure Neumann nullspace is removed by the zero-mean restriction.
\end{proof}

The local Hessian structure of each model is needed for the transfer estimates of
\Cref{sec:c5-hessian}. The bilinear forms at the exact solution are
\begin{align*}
 a_{\rm lin}(v,w)&=\int_\Omega(\nabla v\cdot\nabla w+vw)\dd x,\\
 a_{\rm cub,*}(v,w)&=\int_\Omega(\nabla v\cdot\nabla w+3(u^*)^2vw)\dd x,\\
 a_{\sinh,*}(v,w)&=\int_\Omega(\nabla v\cdot\nabla w+\cosh(u^*)vw)\dd x.
\end{align*}
For a \(p\)-diffusion density, the Hessian matrix in gradient variables is
\[
 \abs\xi^{p-2}I+(p-2)\abs\xi^{p-4}\xi\otimes\xi
\]
at nonzero \(\xi\) in the pure case, with \(\abs\xi^2\) replaced by
\(\eps^2+\abs\xi^2\) in the regularized case. The local Hessian statements that follow concern
\(p\geq2\).

\begin{proposition}[Sufficient continuous local Hessian conditions]
\label{prop:model-hessian}
In dimensions \(d\leq2\), the cubic Hessian is locally Lipschitz as a bilinear form on \(H^1\), and it
is coercive at any solution \(u^*\not\equiv0\). The sinh Hessian is coercive and is locally Lipschitz on
an \(L^\infty\)-bounded neighborhood, with the state perturbation measured in \(H^1\). For fixed
\(\eps>0\), the regularized \(p\)-Hessian is uniformly elliptic, with lower bound proportional to
\(\eps^{p-2}\); its local Lipschitz continuity is valid in the bounded-gradient topology
\(W^{1,\infty}\), not in the \(H^1\) topology asserted for the cubic model.
\end{proposition}

\begin{proof}[Proof sketch]
In dimensions $d\leq2$, Sobolev embedding controls the cubic coefficient on bounded $H^1$ neighborhoods, while an $L^\infty$ bound controls the sinh coefficient.  Differentiating the regularized radial flux gives uniform ellipticity and bounded derivative on bounded-gradient sets.
\end{proof}

\begin{proposition}[Restricted stability for the regularized \(p\)-energy]
\label{prop:regularized-restricted}
Let \(Y_M\subset W^{1,\infty}(\Omega)\) be finite dimensional and let \(\norm\cdot_{0,M}\) be any norm
on \(Y_M\). For fixed \(\eps>0\), bounded subsets of \(Y_M\) have finite restricted constants
\(L_M\) on the corresponding finite-dimensional segments. If the reference norm is chosen so that the gradient seminorm is a
norm on the feasible space, then \(\mu_M>0\). The resulting constants may depend on \(M\), the inverse
inequality on \(Y_M\), and \(\eps\).
\end{proposition}

\begin{proof}[Proof sketch]
On a bounded subset of the finite-dimensional space, norm equivalence and compactness give a finite Hessian-Lipschitz constant.  The positive ellipticity of the regularized density then yields a positive restricted coercivity constant in any reference norm that is a norm on the feasible space.
\end{proof}

The frozen Hessian is constant for the linear model, locally stable for the cubic and sinh models under the stated neighborhood restrictions, uniformly elliptic for fixed $\eps>0$ on bounded-gradient sets, and potentially degenerate for the pure and reaction $p$-models.  The transfer result therefore uses either the local Hessian hypotheses or the global $p$-geometry, according to the model.


Energy coercivity and frozen-Hessian coercivity are distinct. A positive reaction term does
not imply that the frozen \(p\)-Hessian is uniformly positive at points where both the solution and its
gradient vanish, and the continuous degeneracy of the pure model does not preclude the restricted
coercivity in the corresponding finite-dimensional metric.
